% ChargeShield-FL --- DSN 2027 --- Related Work
% Struttura ripresa da ChargeShield-FL_paper_riscrittura_completa_IT.md, S2
% (struttura invariata rispetto alla versione corrente, giudicata buona),
% con S2.3 (auditing empirico della DP) spostata qui da Architecture per
% richiesta esplicita della bozza, e tre nuove citazioni verificate via
% web search il 2026-09-15 (FedMIA, ML-Leaks, Secure Aggregation).

\section{Related Work}
\label{sec:related-work}

\subsection{Membership Inference}
\label{sec:related-mia}

Membership inference against machine learning models was introduced by
Shokri~et~al.~\cite{shokri2017membership}, who trained shadow models to
mimic a target classifier's behavior and used their outputs to train an
attack model that distinguishes members of the training set from
non-members. Yeom~et~al.~\cite{yeom2018privacy} showed that a much
simpler loss-thresholding attack achieves comparable success whenever a
model overfits, tying membership leakage directly to the
generalization gap; our Yeom attack (\texttt{run\_yeom()},
Section~\ref{sec:framework-attacks}) instantiates this construction
unmodified. Salem~et~al.~\cite{salem2019mlleaks} relaxed Shokri's
requirement of knowing the target architecture and training
distribution, showing that a single, architecture-agnostic shadow
model transfers across models and datasets; our Shadow attack
(\texttt{run\_shadow()}) draws on this data/model-independence
argument to justify a single calibrated shadow model per site rather
than an ensemble matched to the target. Carlini~et~al.'s Likelihood
Ratio Attack (LiRA)~\cite{carlini2022membership} reframed the problem
as a per-example hypothesis test between two Gaussian loss
distributions (in, out) and, crucially, argued that the
average-case metrics reported by most prior work --- accuracy, AUC ---
obscure an attack's true capability, which only shows up in the
true-positive rate at low, operationally realistic false-positive
rates. We adopt this argument directly: Section~\ref{sec:results}
reports TPR at fixed FPR (0.1\%, 1\%, 5\%) alongside AUC for exactly
this reason, and our LiRA implementation
(Section~\ref{sec:framework-lira}) is a direct adaptation of Carlini's
construction to the federated, cross-round setting.

A methodologically important and, in our reading, still under-cited
paper is Sablayrolles~et~al.~\cite{sablayrolles2019whitebox}, who
showed that the Bayes-optimal membership inference strategy under a
broad class of threat models is a function of the loss alone, and that
white-box access to intermediate gradients or activations provides no
asymptotic advantage over observing the loss when the adversary's
distributional assumptions are correctly specified. This result
motivates one of the three per-round scoring variants we implement in
\texttt{run\_lira()} (Section~\ref{sec:framework-lira}) alongside the
raw-loss and log-transformed scorers, and it is the reason we treat
score-function choice, not access level, as the primary axis of attack
strength in this paper.

\paragraph{A naming disambiguation.}
A recent and unrelated line of work publishes under the same name as a
module in our own codebase. Zhu~et~al.'s
FedMIA~\cite{zhu2025fedmia} is a membership inference attack that
follows an ``all for one'' principle, aggregating signal from
\emph{every} client across \emph{multiple communication rounds} to
attack a specific target record. Our codebase contains a historical,
never-activated class also named \texttt{FedMIA} (dating to an early
development sprint, explicitly banner-marked as inactive in the source
and never registered in the attack registry) and a distinct, opt-in
diagnostic attack we call FedMIA-gradient
(Section~\ref{sec:framework-attacks}), which reuses two of that
class's numerical routines but implements a different, per-round,
gradient-space reconstruction-error test. Neither of our two
FedMIA-named artifacts is an implementation of, a comparison against,
or a reuse of Zhu~et~al.'s method; the name collision is coincidental
and predates our awareness of their publication. We flag this
explicitly to avoid the double reading a reviewer could otherwise give
our results.

\subsection{Differential Privacy for Federated Learning}
\label{sec:related-dpfl}

The differential privacy framework itself is due to
Dwork~et~al.~\cite{dwork2006calibrating}. Abadi~et~al.'s
DP-SGD~\cite{abadi2016deep} made per-example gradient clipping and
Gaussian noise addition practical at the scale of deep learning
training, and McMahan~et~al. adapted the same mechanism to the
federated, client-level setting first for language
models~\cite{mcmahan2018learning} and, together with the foundational
FedAvg algorithm~\cite{mcmahan2017communication}, established
client-level DP-FedAvg as the standard baseline we compare against
under our \texttt{dp-fedavg} placement (Section~\ref{sec:framework-dp}).
Li~et~al.'s FedProx~\cite{li2020federated} generalizes FedAvg with a
proximal term that reduces to plain FedAvg when its coefficient is
zero, which is how our aggregator is configured for the results in
this paper; we did not exercise heterogeneity-driven proximal
regularization.

A separate line of work asks whether the \emph{aggregation} step
itself, rather than the model updates, can be protected
cryptographically. Bonawitz~et~al.'s secure aggregation
protocol~\cite{bonawitz2017practical} lets a server learn only the sum
of client updates, never an individual client's contribution, using
pairwise masking with dropout resilience. We deliberately do not
assume secure aggregation anywhere in this paper: our threat model
(Section~\ref{sec:threat-model}) grants the aggregator plaintext
access to each client's update, which is what makes the client-level
observation surface in Section~\ref{sec:results-surface-b} attackable
in the first place. Secure aggregation and differential privacy are
complementary, not substitutable, defenses --- the former hides
\emph{who} contributed which update, the latter bounds \emph{how much}
any single record's presence can move the aggregate --- and a
deployment combining both is future work we do not evaluate here.

\subsection{Empirical Auditing of Differential Privacy}
\label{sec:related-auditing}

A formal $(\varepsilon,\delta)$ guarantee bounds worst-case leakage
over every possible adversary and every pair of neighboring datasets;
it does not by itself say how much a specific, instantiated adversary
extracts from a specific trained model. Closing that gap is the
subject of a body of work on empirical DP auditing, and it is the
literature our third ML Plane consumer
(Section~\ref{sec:framework-attacks}) is positioned against, not the
Privacy Auditor's in-line budget accounting
(Section~\ref{sec:framework-auditor}) --- a distinction we found
necessary to state explicitly after an earlier draft of this paper
conflated the two (Section~\ref{sec:framework-auditor}).
Jagielski~et~al.~\cite{jagielski2020auditing} and
Nasr~et~al.~\cite{nasr2021adversary} developed adversary-instantiation
techniques that turn a DP mechanism's theoretical bound into an
empirical lower bound on $\varepsilon$ by constructing worst-case
canary inputs and measuring how reliably an attacker distinguishes
their inclusion. Steinke~et~al.~\cite{steinke2023privacy} showed that
a useful empirical lower bound can be obtained from a \emph{single}
training run rather than the many thousands classical auditing
requires, by injecting multiple independent canaries and analyzing
them jointly --- a result directly relevant to any future single-run
canary campaign on this codebase. Maddock~et~al.'s
CANIFE~\cite{maddock2023canife} adapts canary-based auditing
specifically to the federated setting, addressing the practical
difficulty that in cross-silo FL a canary must be crafted per client
and per round rather than once for a centralized dataset; our own
canary positive control (Section~\ref{sec:framework-canary}) is a
direct, simplified descendant of this idea, restricted to verifying
attack sensitivity rather than deriving a formal $\varepsilon$ lower
bound.
Jayaraman and Evans~\cite{jayaraman2019evaluating} is the paper whose
framing we adopt for reading a null attack result against a DP
mechanism at all: they show empirically that the $\varepsilon$ values
DP-SGD implementations use in practice are far looser than the
$\varepsilon$ at which membership inference actually degrades, so a
null result at a \emph{practical} $\varepsilon$ is uninformative about
whether DP caused it unless the same attack is also shown to fail
without any DP at all. Section~\ref{sec:framework-dp}'s
\texttt{dp-fedavg} upper-bound placement and
Section~\ref{sec:results}'s no-DP baseline exist specifically to make
that comparison available. Rahman~et~al.~\cite{rahman2018membership}
is the closest prior empirical study to ours in spirit --- membership
inference against differentially private deep learning models on
real, non-benchmark data --- though restricted to the centralized
setting; we return to their findings when discussing the practical
$\varepsilon$--utility trade-off in Section~\ref{sec:discussion}.

\subsection{Byzantine-Robust and Adjacent Federated Learning}
\label{sec:related-byzantine}

Blanchard~et~al.'s Krum~\cite{blanchard2017machine} and the broader
literature on Byzantine-tolerant aggregation address a threat model
orthogonal to the one this paper measures: a client that deviates from
the protocol to corrupt the global model, versus a server that follows
the protocol faithfully but observes passively. Our ByzantineDetector
component (Section~\ref{sec:framework-byzantine}) draws on this
literature for its Krum and cosine-similarity mechanisms, but --- as
we state explicitly in Section~\ref{sec:framework-byzantine} after
finding and correcting an earlier conflation of the two threat models
in our own writing --- no integrity mechanism from this literature
provides any confidentiality guarantee against the honest-but-curious
aggregator this paper's results characterize, and no privacy metric in
this paper detects Byzantine behavior. Bagdasaryan~et~al.'s backdoor
attack~\cite{bagdasaryan2020backdoor} is a further example of a
client-side threat in this same orthogonal category. Kairouz~et~al.'s
survey~\cite{kairouz2021advances} and Truong~et~al.'s GDPR-oriented
survey~\cite{truong2021privacy} both situate privacy and integrity as
distinct concerns within the broader federated learning literature,
consistent with how we separate them architecturally
(Section~\ref{sec:framework}).

\subsection{Relationship to Our Prior Work}
\label{sec:related-prior}

This paper extends our own prior work~\cite{imperatrice2026toward}, a
QRS 2026 fast abstract that first proposed the ML Plane as an
architectural response to the Purdue-model observability gap FL
introduces (Section~\ref{sec:introduction}). That prior work proposed
the substrate without a registered consumer: no in-line budget
auditor, no integrity detector, and no offline attack suite were
implemented or evaluated against it. The contribution of the present
paper is to instrument that substrate with three independent
consumers, deploy it against a real, multi-site EV charging dataset,
and report what each consumer actually measures --- including, in
Section~\ref{sec:validation}, a methodology for telling a genuine null
membership-inference result apart from an insensitive one, which did
not exist in any form in the prior abstract.

\emph{Editorial note preserved from the existing bibliography:} DSN's
double-blind anonymization rules require self-citation in the third
person without omission (see CFP, Anonymization Rules); the wording
above is written to satisfy that requirement and should not be altered
to first person or removed prior to submission.
